%
% notification-race.tex
%
% Isolated verification of the pendingNotify mutex (Section~\ref{sec:race}
% of notification.tex).  The full NotificationState carries thirty-odd
% components for the talk/wall/read marker channels, none of which
% participate in the race; reusing it for this proof made the reachable
% space explode combinatorially (over 6.5M states and still growing when
% killed) without adding anything the race property depends on.  This
% file lifts out exactly the five components the race and the mutex
% actually touch --- $pendingNotify$, $inFlight$, the two per-site
% snapshots, and the $raceLost$ witness --- as their own tiny machine,
% checkable exhaustively in milliseconds.  It is cited from
% notification.tex \S~\emph{The Race: Three Critical Sections, One Shared Flag} rather than substituting for it: the
% schemas here are verbatim restrictions of the ones there, and the two
% files are read together, not in place of one another.
%

\documentclass[a4paper,10pt,fleqn]{article}
\usepackage[margin=1in]{geometry}
\usepackage{fuzz}

\begin{document}

\title{The \texttt{pendingNotify} Mutex: An Isolated Proof}
\author{Auxiliary to \emph{Biff Notification System: A Z Specification}}
\date{September 2026}

\maketitle

\section{Purpose}

\texttt{notification.tex} \S~\emph{The Race: Three Critical Sections, One Shared Flag} models three critical sections
--- suspenders, belt, capture/reconnect --- each split into a read/decide
step and a write step, and proves a mutex serialising them closes a
lost-update race on $pendingNotify$.  That proof is stated there against
the full $NotificationState$ schema, which also carries the talk and
wall marker machinery (Sections~2--10 of the main document) so that the
race schemas can share $NotificationState$'s existing framing
conventions.  Combined, the reachable space is large enough that an
exhaustive check --- even at the floor cardinality every carrier in the
main document already uses --- did not terminate in a reasonable time or
memory budget.  None of that machinery participates in the race: the
race is a five-variable property.  This file states it as its own
machine and checks it exhaustively.

\section{State}

\begin{zed}
  ZBOOL ::= ztrue | zfalse
\end{zed}

\begin{zed}
  RaceSite ::= raceSuspenders | raceBelt | raceCapture
\end{zed}

\begin{schema}{RaceState}
  pendingNotify : ZBOOL \\
  inFlight : \power RaceSite \\
  beltSnapshot : ZBOOL \\
  capSnapshot : ZBOOL \\
  raceLost : ZBOOL \\
  notifyLost : ZBOOL
  \where
  \# inFlight \leq 1 \\
  raceLost = zfalse \\
  notifyLost = zfalse
\end{schema}

$pendingNotify$ is the shared flag every site reads and writes.
$inFlight$ tracks which sites have completed their read/decide step but
not yet their write --- the gap a non-atomic \texttt{await} opens up in
the real code.  $beltSnapshot$ and $capSnapshot$ are what each
``clearing'' site read at the moment it entered that gap;
$raceSuspenders$ never clears (its only write is an unconditional
re-arm, \emph{Suspenders Failure and the No-Recovery Sink} of the main document), so it needs
no snapshot.  $raceLost$ is round~2's witness: it fires when a clearing
site's eventual write is based on a value that has since changed
underneath it.  $notifyLost$ is round~1's witness, restated at the
current three-site granularity: it fires when a site's failure commit
leaves $pendingNotify$ un-re-armed.  The invariant $\# inFlight \leq 1$
is the mutex; $raceLost = zfalse$ and $notifyLost = zfalse$ are the two
safety properties it and the round~1 fix are meant to establish jointly.
All are stated here exactly as they are in \texttt{notification.tex},
restricted to the variables that matter.

\section{Initialisation}

\begin{schema}{RaceInit}
  RaceState'
  \where
  pendingNotify' = zfalse \\
  inFlight' = \emptyset \\
  beltSnapshot' = zfalse \\
  capSnapshot' = zfalse \\
  raceLost' = zfalse \\
  notifyLost' = zfalse
\end{schema}

\section{Suspenders (set-only)}

The suspenders site's only effect on $pendingNotify$ is an unconditional
re-arm on failure (\texttt{notification.tex}'s $PollTickNotifyFail$); its
success commit does not touch $pendingNotify$ at all and is therefore
omitted here --- it contributes nothing to reachability of $raceLost$ or
$notifyLost$.  $SusCommitSetUnrecorded$ is round~1's defect at this
site, restated: the commit fires but leaves $pendingNotify$ untouched
instead of re-arming it.

\begin{schema}{SusBegin}
  \Delta RaceState
  \where
  raceSuspenders \notin inFlight \\
  inFlight' = inFlight \cup \{ raceSuspenders \} \\
  pendingNotify' = pendingNotify \\
  beltSnapshot' = beltSnapshot \\
  capSnapshot' = capSnapshot \\
  raceLost' = raceLost \\
  notifyLost' = notifyLost
\end{schema}

\begin{schema}{SusCommitSet}
  \Delta RaceState
  \where
  raceSuspenders \in inFlight \\
  inFlight' = inFlight \setminus \{ raceSuspenders \} \\
  pendingNotify' = ztrue \\
  beltSnapshot' = beltSnapshot \\
  capSnapshot' = capSnapshot \\
  raceLost' = raceLost \\
  notifyLost' = notifyLost
\end{schema}

\begin{schema}{SusCommitSetUnrecorded}
  \Delta RaceState
  \where
  raceSuspenders \in inFlight \\
  inFlight' = inFlight \setminus \{ raceSuspenders \} \\
  pendingNotify' = pendingNotify \\
  beltSnapshot' = beltSnapshot \\
  capSnapshot' = capSnapshot \\
  raceLost' = raceLost \\
  notifyLost' = ztrue
\end{schema}

\section{Belt}

\begin{schema}{BeltBegin}
  \Delta RaceState
  \where
  raceBelt \notin inFlight \\
  inFlight' = inFlight \cup \{ raceBelt \} \\
  beltSnapshot' = pendingNotify \\
  pendingNotify' = pendingNotify \\
  capSnapshot' = capSnapshot \\
  raceLost' = raceLost \\
  notifyLost' = notifyLost
\end{schema}

\begin{schema}{BeltCommitClear}
  \Delta RaceState
  \where
  raceBelt \in inFlight \\
  inFlight' = inFlight \setminus \{ raceBelt \} \\
  pendingNotify' = zfalse \\
  ((pendingNotify \neq beltSnapshot) \implies raceLost' = ztrue) \\
  ((pendingNotify = beltSnapshot) \implies raceLost' = raceLost) \\
  beltSnapshot' = beltSnapshot \\
  capSnapshot' = capSnapshot \\
  notifyLost' = notifyLost
\end{schema}

\begin{schema}{BeltCommitSet}
  \Delta RaceState
  \where
  raceBelt \in inFlight \\
  inFlight' = inFlight \setminus \{ raceBelt \} \\
  pendingNotify' = ztrue \\
  beltSnapshot' = beltSnapshot \\
  capSnapshot' = capSnapshot \\
  raceLost' = raceLost \\
  notifyLost' = notifyLost
\end{schema}

\begin{schema}{BeltCommitSetUnrecorded}
  \Delta RaceState
  \where
  raceBelt \in inFlight \\
  inFlight' = inFlight \setminus \{ raceBelt \} \\
  pendingNotify' = pendingNotify \\
  beltSnapshot' = beltSnapshot \\
  capSnapshot' = capSnapshot \\
  raceLost' = raceLost \\
  notifyLost' = ztrue
\end{schema}

\section{Capture/Reconnect Flush}

$CapBegin$'s guard $pendingNotify = ztrue$ mirrors
$CaptureSessionFlushBegin$'s in the main document: capture only enters
this critical section when there is something recorded to flush.

\begin{schema}{CapBegin}
  \Delta RaceState
  \where
  raceCapture \notin inFlight \\
  pendingNotify = ztrue \\
  inFlight' = inFlight \cup \{ raceCapture \} \\
  capSnapshot' = pendingNotify \\
  pendingNotify' = pendingNotify \\
  beltSnapshot' = beltSnapshot \\
  raceLost' = raceLost \\
  notifyLost' = notifyLost
\end{schema}

\begin{schema}{CapCommitClear}
  \Delta RaceState
  \where
  raceCapture \in inFlight \\
  inFlight' = inFlight \setminus \{ raceCapture \} \\
  pendingNotify' = zfalse \\
  ((pendingNotify \neq capSnapshot) \implies raceLost' = ztrue) \\
  ((pendingNotify = capSnapshot) \implies raceLost' = raceLost) \\
  beltSnapshot' = beltSnapshot \\
  capSnapshot' = capSnapshot \\
  notifyLost' = notifyLost
\end{schema}

\begin{schema}{CapCommitSet}
  \Delta RaceState
  \where
  raceCapture \in inFlight \\
  inFlight' = inFlight \setminus \{ raceCapture \} \\
  pendingNotify' = ztrue \\
  beltSnapshot' = beltSnapshot \\
  capSnapshot' = capSnapshot \\
  raceLost' = raceLost \\
  notifyLost' = notifyLost
\end{schema}

\begin{schema}{CapCommitSetUnrecorded}
  \Delta RaceState
  \where
  raceCapture \in inFlight \\
  inFlight' = inFlight \setminus \{ raceCapture \} \\
  pendingNotify' = pendingNotify \\
  beltSnapshot' = beltSnapshot \\
  capSnapshot' = capSnapshot \\
  raceLost' = raceLost \\
  notifyLost' = ztrue
\end{schema}

\section{Verification}

Eleven operations, six state components, $\# inFlight \leq 1$ bounding
$inFlight$ to four possible values: the whole reachable space is small
by construction.  Two runs, exactly as in \texttt{notification.tex}
\S~\emph{The Race: Three Critical Sections, One Shared Flag}, at $DEFAULT\_SETSIZE~2$ (there are no carrier-typed
components left to scope, so the setsize binds nothing here --- it is
given only for invocation parity with the main document):

\paragraph{Unsafe: the race, reproduced.}  All three invariant conjuncts
removed from $RaceState$ (an operation schema's implicit
$RaceState \land RaceState'$ framing means \emph{all} must go---leaving
$raceLost = zfalse$ or $notifyLost = zfalse$ in place while dropping
only the mutex would still forbid every transition that sets either
witness to $ztrue$, hiding the defect instead of exposing it, exactly
as $notifyLost = zfalse$ had to be dropped whole in
\texttt{notification.tex} \S~\emph{Suspenders Failure and the
No-Recovery Sink}):

\begin{verbatim}
probcli docs/notification-race.tex -model_check -noltl -coverage \
  -goal "raceLost = ztrue" \
  -p DEFAULT_SETSIZE 2 -p MAX_OPERATIONS 2000
\end{verbatim}

\begin{verbatim}
Found GOAL in state id 15
States analysed: 30   Transitions fired: 129
*** COUNTER EXAMPLE FOUND ***
*** TRACE (length=5):
 1: INITIALISATION(beltSnapshot=zfalse,capSnapshot=zfalse,inFlight={},
    notifyLost=zfalse,pendingNotify=zfalse,raceLost=zfalse)
 2: SusBegin
 3: BeltBegin
 4: SusCommitSet
 5: BeltCommitClear(inFlight={},pendingNotify=zfalse,raceLost=ztrue)
\end{verbatim}

This is precisely the belt/suspenders interleaving described informally
in \texttt{notification.tex} \S~\emph{The Race: Three Critical Sections, One Shared Flag}: both sites read/decide
before either writes (steps~2--3; belt's read of $pendingNotify =
zfalse$ decides ``clear''); suspenders commits first,
$pendingNotify' = ztrue$ (step~4), \emph{while belt's write is still
outstanding}; belt's commit (step~5) then writes $zfalse$ over it,
using the decision it made before the re-arm existed.  $pendingNotify$
ends $zfalse$ and $raceLost$ records that the write was stale when it
landed.

The same file, same goal search technique, reproduces round~1's defect
at the new three-site granularity in two steps:

\begin{verbatim}
probcli docs/notification-race.tex -model_check -noltl -coverage \
  -goal "notifyLost = ztrue" \
  -p DEFAULT_SETSIZE 2 -p MAX_OPERATIONS 2000
\end{verbatim}

\begin{verbatim}
Found GOAL in state id 5
States analysed: 2   Transitions fired: 9
*** COUNTER EXAMPLE FOUND ***
*** TRACE (length=3):
 1: INITIALISATION(beltSnapshot=zfalse,capSnapshot=zfalse,inFlight={},
    notifyLost=zfalse,pendingNotify=zfalse,raceLost=zfalse)
 2: BeltBegin
 3: BeltCommitSetUnrecorded
\end{verbatim}

A belt send fails, and the commit that should re-arm $pendingNotify$
leaves it untouched instead: two steps, no interleaving required---this
is round~1's defect, not round~2's, restated for belt rather than
suspenders and reachable on its own.

\paragraph{Safe: the mutex plus the re-arm, exhaustive.}  $RaceState$ as
committed above ($\# inFlight \leq 1$, $raceLost = zfalse$, and
$notifyLost = zfalse$ all present):

\begin{verbatim}
probcli docs/notification-race.tex -model_check -noltl -coverage \
  -p DEFAULT_SETSIZE 2 -p MAX_OPERATIONS 2000
\end{verbatim}

\begin{verbatim}
Model checking finished, all open states visited
Model checking time: 33 ms   States analysed: 22   Transitions fired: 41
No counter example found. ALL states visited.
Coverage: [STATES (23), deadlocked:0, invariant_violated:0,
  invariant_not_checked:0, open:0, COVERED_OPERATIONS (9): BeltBegin,
  BeltCommitClear, BeltCommitSet, CapBegin, CapCommitClear,
  CapCommitSet, INITIALISATION, SusBegin, SusCommitSet,
  UNCOVERED_OPERATIONS (3): BeltCommitSetUnrecorded,
  CapCommitSetUnrecorded, SusCommitSetUnrecorded]
\end{verbatim}

Twenty-three states, complete, in 33~ms.  The three ``Unrecorded''
defect operations---round~1's finding, now stated and proved dead at
all three sites in one pass---are reported uncovered, the same proof
shape as $WithdrawStale$ in the main document.  $raceLost = zfalse$ and
$notifyLost = zfalse$ both hold across the entire reachable space.
These results are cited from \texttt{notification.tex} \S~\emph{The Race: Three Critical Sections, One Shared Flag}
alongside that document's own full-model goal-directed counterexample
(the same interleaving, found via a different pair of colliding sites
because ProB's BFS explores the larger model in a different order);
this file exists so the \emph{exhaustive} half of the claim rests on a
check that terminates in milliseconds, not on the full model's
combinatorial product with unrelated marker machinery, which did not
terminate in a reasonable time or memory budget when tried directly.

\end{document}
